\chapter{如何阅读一条编年事件}

\section{三层叙事}

每个时期先给出制度与地缘背景，然后按照年代推进，最后在大转折处暂停并深描。读者可以顺着时间线阅读，也可以通过因果链回看某项改革、战争、继承危机或财政制度如何积累影响。

\section{条目格式}

\begin{event}{公元前841年}{周厉王出奔}{可信年表}{镐京及国人政治}
\term{经过}：国人起而逐厉王，周公、召公共同执政，传统纪年称为“共和”。

\term{前因}：王室对山林川泽等资源的控制与政治排斥，引发贵族和国人的冲突；具体政策细节主要来自后出的叙事，不能把它们全部当作同时代行政档案。

\term{后果}：共和元年成为传统连续编年的起点；更重要的是，事件显示西周王室调停贵族与国人利益的能力已显著削弱。
\end{event}

\section{因果关系的标签}

\begin{tabularx}{\linewidth}{@{}l X@{}}
\toprule
标签 & 含义 \\
\midrule
结构性原因 & 长期存在的制度、资源、人口或地缘压力 \\
触发因素 & 使既有矛盾在某一时间点爆发的直接事件 \\
使能条件 & 让行动者能采取某种行动的联盟、技术或信息条件 \\
直接后果 & 数月到数年内可观察的政权、政策或战争结果 \\
中长期影响 & 改变后续选择、制度路径或地区力量结构的结果 \\
反馈回路 & 某个后果又反过来加剧原先矛盾 \\
\bottomrule
\end{tabularx}

\begin{sourcenote}
本书优先使用传世史籍、出土文献与现代研究的交叉信息。对西周早期、战国纪年、人物动机等争议问题，不会用单一叙事抹平分歧。
\end{sourcenote}
